\section{本章小结}
\label{sec:lpipm-summary}

本章开头问：能不能不贴着边走，从可行域的内部穿过去？答案是把互补松弛$x_is_i=0$放宽成$x_is_i=\mu$。原始和对偶都有严格可行点时，这组方程对每个$\mu>0$有唯一解，解连成中心路径，沿路径的对偶间隙恰好是$n\mu$。路径的极限是一对严格互补的最优解，落在最优面的解析中心，解不唯一时并不是顶点。

沿路径前进靠Newton法。线性化$x_is_i=\sigma\mu$得到一个由原始残差、对偶残差和互补残差驱动的线性方程组，消元后只需解$m\times m$的对称正定方程$\bm{A}\bm{D}\bm{A}^\top\Delta\bm{y}=\bm{h}$。到边界的比例规则保持$\bm{x},\bm{s}>\bm{0}$，$\sigma$在减小互补量和靠近中心之间折中，预测—校正用一次试算选定$\sigma$并补偿二阶项，两次求解共用一次分解。

这些步骤合起来是一个完整的求解过程。起点可以严格可行，也可以只要求正性，让等式残差和互补量一起下降；需要判别不可行或无界时，用预处理和齐次自对偶嵌入。停止时要同时看三项残差，只有在原始、对偶都可行时，$\bm{x}^\top\bm{s}$才是目标误差的上界，否则还要计入等式残差被最优解放大的部分。

交付的结果有几个层次，它们的含义不同：浮点容差下的近似解，基恢复之后的基，用精确算术核对过的证书。理论方面，短步路径跟踪的迭代次数不超过$O(\sqrt{n}\log(n\mu^{(0)}/\varepsilon))$，配合起点构造和位长分析得到线性规划的多项式算法；椭球法更早建立了这一可解性。理论规则和实用的预测—校正之间有明显差距，不能把后者的经验速度当作定理。

选择算法时，单纯形法便于热启动、返回基，适合反复求解和灵敏度分析；内点法迭代次数少，适合大规模稀疏问题，但对$\bm{A}\bm{D}\bm{A}^\top$的填充敏感，需要基时要另作恢复。两者都需要在最后检验可行性和最优性。

路径跟踪的思路并不限于线性规划。当约束是一般的光滑函数或锥时，障碍和扰动互补仍然可用，但方程不再是双线性的，Newton系统含有Hessian和Jacobian的变化，严格内点、初值和停止证书都有新的条件，这些由第~\ref{chap:nonlinear-interior-point}章展开。
